\documentclass[10pt,a4paper]{amsart}
\usepackage[margin=19mm]{geometry}
\usepackage{amsmath,amssymb,mathtools}
\usepackage[scr=rsfs,cal=euler]{mathalfa}
\usepackage{xfrac}
\usepackage[unicode,hidelinks,bookmarks=false]{hyperref}
\newcommand{\ie}{i.e.\ }
\newcommand{\cf}{cf.\ }
\newcommand{\resp}{respectively\ }
\newcommand{\Subsection}{\subsection}
\providecommand{\nobreakdash}{\nobreak\hbox{-}}
\setlength{\emergencystretch}{2em}
\multlinegap0pt
\usepackage{bm}
\usepackage{bbm}
\usepackage{dsfont}
\usepackage{xspace}
\usepackage{cancel}
\usepackage{mathtools}
\usepackage{amsthm}
\makeatletter
\@ifundefined{thm}{\newtheorem{thm}{Theorem}}{}
\makeatother
\newcommand{\Soc}{\operatorname{Soc} }
\title[Skew brace extensions, second cohomology and complements]{Skew brace extensions, second cohomology and complements}
\author{Nishant Rathee, Manoj K. Yadav}
\date{Source: arXiv:2601.12371v2 (CC BY 4.0)}
\begin{document}
\maketitle

\noindent\emph{Source and licence.} Skew brace extensions, second cohomology and complements. Version arXiv:2601.12371v2 (CC BY 4.0), distributed under the Creative Commons Attribution 4.0 International licence, as recorded in the arXiv licence metadata for this exact version and shown on the abstract page https://arxiv.org/abs/2601.12371v2. Full rights notice: licenses/arxiv-2601.12371-CC-BY-4.0.txt.\par\smallskip

\noindent\textit{Editorial scope.} Counted unit: the Thm whose source label is \texttt{prop-split-ext} in the article (source0.tex lines 445--447, source label \texttt{prop-split-ext}), delivered together with the article's own complete original proof of it (source0.tex lines 449--510, one \texttt{proof} environment of 62 lines). No mathematics is written, repaired or re-derived: statement and proof are the article's own text line for line. Required context retained uncounted: sb-ext (lines 267--269); g-ext (lines 271--273). Every definition, display and label of those lines is retained unchanged. Omitted from this package and not redistributed: 1--266, 270--270, 274--444, 448--448, 511--1776. Nothing inside a retained excerpt is omitted or altered: every retained body is delivered exactly as the article's lines are written, so no omission receipt of any kind accompanies this package. Citations: the retained excerpts carry no citation command at all, so no bibliography entry of the article is transcribed and no reference is redistributed. Reference adaptation: none was needed and none was made; every reference inside the delivered unit resolves inside it, so no unresolved reference or citation is produced. Printed slips kept literal: no printed word, symbol, hypothesis or number of the counted statement or of its proof was altered. Layout and class: the article's own class is \texttt{amsart} with packages including amsmath, amsfonts, amssymb, upgreek, amscd, mathtools, hyperref, tikz, enumerate, tikz-cd, graphicx; the wrapper is the lane's pinned \texttt{amsart} editorial wrapper at 10pt on A4, which restates the theorem-like environments the retained text needs and adds the article's own macro definitions that the retained text needs (\texttt{\textbackslash Soc}), with extra wrapper packages (bm, bbm, dsfont, xspace, cancel, mathtools). The delivered numbering is the unit's own. Assets: no retained span contains a figure environment, a graphics-inclusion command, a PDF-page inclusion, a table or a TikZ picture; no image file is copied into this package, the package declares no asset dependency and no image file is redistributed with it. Licence: version arXiv:2601.12371v2 of this article is distributed under the Creative Commons Attribution 4.0 International licence, as recorded in the arXiv licence metadata for this exact version and shown on the abstract page https://arxiv.org/abs/2601.12371v2. A full rights notice is delivered at licenses/arxiv-2601.12371-CC-BY-4.0.txt. Characters: every retained excerpt is pure ASCII, so the delivered engine, XeLaTeX with its default Latin Modern text font, reports no missing character.
\begin{equation}\label{sb-ext}
0 \longrightarrow I \xrightarrow{i} E \xrightarrow{\pi} H \longrightarrow 0,
\end{equation}

\begin{equation}\label{g-ext}
0 \longrightarrow \Lambda_I  \xrightarrow{i \times i} \Lambda_E \xrightarrow{\pi \times \pi} \Lambda_H \longrightarrow 0.
\end{equation}

\begin{thm}\label{prop-split-ext}
	The skew brace extension \eqref{sb-ext}  with $I \le \Soc(E)$ splits if and only if  the associated group extension \eqref{g-ext} splits.
\end{thm}

\begin{proof}
Note that `only if' part of the assertion follows from the construction.  So assume that the group extension \eqref{g-ext} splits. Then there exists a section (group homomorphism) 
$s : \Lambda_H \to \Lambda_E$ such that  \( (\pi \times \pi) \circ s = \mathrm{Id}_{\Lambda_H} \). For all \( h, g \in H \), we can write
$$
s(h, g) = (s_1(h, g), s_2(h, g)),$$
where \( s_1, s_2 : H \times H \to E \) are the component maps of $s$.
Since \( s \) is a group homomorphism, for all \( (h_1, g_1), (h_2, g_2) \in \Lambda_H \), we have
$$s(h_1 + \lambda_{g_1}(h_2),\; g_1 \circ g_2) = s(h_1, g_1) \cdot s(h_2, g_2).$$
Now writing in the  component form and making further computations we get
$$s(h_1 + \lambda_{g_1}(h_2),\; g_1 \circ g_2) = (s_1(h_1 + \lambda_{g_1}(h_2),\; g_1 \circ g_2),\; s_2(h_1 + \lambda_{g_1}(h_2),\; g_1 \circ g_2)) $$
and 
\begin{align*}
s(h_1, g_1) \cdot s(h_2, g_2)  &= (s_1(h_1, g_1), s_2(h_1, g_1)) \cdot (s_1(h_2, g_2), s_2(h_2, g_2)) \\
		&= \left(s_1(h_1, g_1) + \lambda_{s_2(h_1, g_1)}(s_1(h_2, g_2)),\; s_2(h_1, g_1) \circ s_2(h_2, g_2)\right).
	\end{align*}
	Comparing the first components, we obtain the key relation:
	\begin{equation}\label{eq:main1}
		s_1(h_1 + \lambda_{g_1}(h_2),\; g_1 \circ g_2) = s_1(h_1, g_1) + \lambda_{s_2(h_1, g_1)}(s_1(h_2, g_2)).
	\end{equation}
	
	Now, $s$ being a section, we see that \( \pi(s_1(h, g)) = h \) and \( \pi(s_2(h, g)) = g \), and hence the difference \( s_2(h, g) - s_1(g, h) \in I \). So there exists a  map \( f : H \times H \to I \) such that 
		\[
	s_2(h, g) = s_1(g, h) + f(h, g).
	\]
 In particular, $s_1(0, h)$ and $s_2(h, 0)$ both lie in $I$ for all $h \in H$. Observe that for any $h_1, h_2 \in H$, we have
	\[
	s(h_1, h_2) = s((h_1, 0)(0, h_2)),
	\]
	which, using \eqref{eq:main1},  implies that
	\begin{equation}\label{eq:addsplit}
		s_1(h_1, h_2) = s_1(h_1, 0) + s_1(0, h_2),
	\end{equation}
	since $s_2(h_1, 0) \in I$.

	
Using the fact that $I \le \Soc(E)$ and the identity $s_2(h, g) = s_1(g, h) + f(h, g)$, \eqref{eq:main1} takes the form
	\begin{equation}\label{eq:main2}
		s_1(h_1 + \lambda_{g_1}(h_2),\; g_1 \circ g_2) = s_1(h_1, g_1) + \lambda_{s_1(g_1, h_1)}(s_1(h_2, g_2)).
	\end{equation}
Define a map \( t : H \to E \) by
	\[
	t(h) := s_1(h, 0).
	\]
	That $t$ is a st-section of the projection \( \pi : E \to H \) is obvious. We claim that \( t \) is a skew brace homomorphism. Taking  \( g_1 = g_2 = 0 \) in  \eqref{eq:main2}, we get
	$$s_1(h_1 + h_2, 0) = s_1(h_1, 0) + \lambda_{s_1(0, h_1)}(s_1(h_2, 0))$$
	for all  \( h_1, h_2 \in H \). 	Since \( s_1(0, h_1) \in I \), this, in the form of $t$, gives
	\[
	t(h_1 + h_2) = t(h_1) + t(h_2).
	\]
This shows that \( t \) is a homomorphism of the additive groups. Next, taking  \( h_1 = 0 \) and \( g_2 = 0 \) in \eqref{eq:main2}, we get
$$s_1(\lambda_{g_1}(h_2), g_1) = s_1(0, g_1) + \lambda_{s_1(g_1, 0)}(s_1(h_2, 0))$$
for all $g_1, h_2 \in H$, which,
	using  \eqref{eq:addsplit} on the left hand side and the fact $I \le \Soc(E)$, gives
	\[
	s_1(\lambda_{g_1}(h_2), 0) = \lambda_{s_1(g_1, 0)}(s_1(h_2, 0)).
	\]
	This, using the definition of $t$, gives 
	\[
	t(\lambda_{g_1}(h_2)) = \lambda_{t(g_1)}(t(h_2)).
	\]
Since $t$ is  a homomorphism of the additive groups, we get $t(g_1 \circ h_2) = t(g_1) \circ t(h_2)$ for all $g_1, h_2 \in H$. Hence $t$ is a skew brace homomorphism, and the proof is complete. 		
\end{proof}

\end{document}
